% from ega4-9-fr.tex
\oldpage[IV-3]{77}

\begin{lemma}[9.7.2]
\label{IV.9.7.2-fr}
Soient $S$ un préschéma irréductible de point générique $\eta$, $f:X\to S$
un morphisme de présentation finie. Si $X_\eta$ n'est pas réduit, il existe un
voisinage ouvert $U$ de $\eta$ dans $S$ tel que, pour tout $s\in U$, $X_s$ ne
soit pas réduit.
\end{lemma}

En effet, soit $\mathcal N$ le Nilradical de $\mathcal O_X$~; il résulte de
(8.2.13), appliqué suivant la méthode de (8.1.2, \emph{a)}) que $\mathcal N_\eta$
est le nilradical de $\mathcal O_{X_\eta}$, et l'hypothèse est que
$\mathcal N_\eta\neq0$. On conclut de (9.4.5) et (9.2.3) qu'il y a un voisinage
$U$ de $\eta$ tel que, pour tout $s\in U$, $\mathcal N_s$ s'identifie à un Idéal
de $\mathcal O_{X_s}$ et $\mathcal N_s\neq0$~; comme $\mathcal N_s$ est
évidemment contenu dans le Nilradical de $\mathcal O_{X_s}$, on voit que $X_s$
n'est pas réduit pour $s\in U$.

\phantomsection\label{IV.9.7.3-fr}%
\textbf{(9.7.3)} Étant donné un polynôme
$F\in A[\mathrm T_1,\ldots,\mathrm T_n]$, où $A$ est un anneau et les
$\mathrm T_i$ des indéterminées, pour tout homomorphisme d'anneaux
$\rho:A\to B$, nous désignerons par $F_{(\rho)}$ ou $F_{(B)}$ le polynôme de
$B[\mathrm T_1,\ldots,\mathrm T_n]$ obtenu en remplaçant chaque coefficient de
$F$ par son image par $\rho$. Si $k$ est un corps,
$F\in k[\mathrm T_1,\ldots,\mathrm T_n]$ un polynôme non constant et
$X=\operatorname{Spec}(k[\mathrm T_1,\ldots,\mathrm T_n]/(F))$, dire que $X$
est intègre (ou que l'idéal $(F)$ est premier) signifie que $F$ est
\emph{irréductible} (c'est-à-dire que dans toute factorisation $F=F_1F_2$ en
polynômes de $k[\mathrm T_1,\ldots,\mathrm T_n]$, $F_1$ ou $F_2$ est de degré
$0$)~; cela résulte de ce que l'anneau $k[\mathrm T_1,\ldots,\mathrm T_n]$ est
factoriel. On en déduit aussitôt ((4.6.2) et (4.5.2))~:

\begin{lemma}[9.7.4]
\label{IV.9.7.4-fr}
Soient $k$ un corps, $\Omega$ une extension algébriquement close de $k$, $F$ un
polynôme non constant de $k[\mathrm T_1,\ldots,\mathrm T_n]$. Les conditions
suivantes sont équivalentes~:
\begin{enumerate}[label=\emph{\alph*)}]
  \item $X=\operatorname{Spec}(k[\mathrm T_1,\ldots,\mathrm T_n]/(F))$ est
  géométriquement intègre.
  \item $F_{(K)}$ est irréductible pour toute extension $K$ de $k$.
  \item $F_{(\Omega)}$ est irréductible.
\end{enumerate}
Nous dirons dans ce cas que $F$ est \emph{géométriquement irréductible}.
\end{lemma}

\begin{lemma}[9.7.5]
\label{IV.9.7.5-fr}
Soient $A$ un anneau intègre, $K$ son corps des fractions, $F$ un polynôme non
constant de $A[\mathrm T_1,\ldots,\mathrm T_n]$ de degré $d$ tel que $F_{(K)}$
soit géométriquement irréductible. Alors il existe $f\neq0$ dans $A$ tel que
pour tout $x\in D(f)$, $F_{(\boldsymbol{k}(x))}$ soit géométriquement
irréductible.
\end{lemma}

Posons $F=\sum_\alpha c_\alpha\mathrm T^\alpha$ avec comme d'ordinaire
$c_\alpha=c_{\alpha_1\alpha_2\ldots\alpha_n}$,
$\mathrm T^\alpha=\mathrm T_1^{\alpha_1}\mathrm T_2^{\alpha_2}\cdots
\mathrm T_n^{\alpha_n}$,
$|\alpha|=\alpha_1+\alpha_2+\cdots+\alpha_n\leqslant d$ (un au moins des
$c_\alpha$ tels que $|\alpha|=d$ n'étant pas nul). Comme les $c_\alpha$ non
nuls sont inversibles dans $K$, on peut supposer, en remplaçant au besoin $A$
par un anneau $A_g$ ($g\neq0$ dans $A$) que les $c_\alpha$ non nuls sont
inversibles dans $A$. Il en résulte que pour tout $x\in\operatorname{Spec}(A)$,
$F_{(\boldsymbol{k}(x))}$ est \emph{de degré $d$}.

Nous démontrerons d'abord un lemme préliminaire.

\begin{lemma}[9.7.5.1]
\label{IV.9.7.5.1-fr}
Soient $A$ un anneau intègre, $S=\operatorname{Spec}(A)$, $\eta$ le point
générique de $S$, $B$ l'anneau de polynômes
$A[\mathrm T_1,\ldots,\mathrm T_m]$, $(P_i)_{i\in I}$ une famille finie
d'éléments de $B$, $\mathfrak a$ l'idéal de $B$ engendré par les $P_i$. Pour
tout $s\in S$, soit $\mathfrak a_s$ l'idéal de
$\boldsymbol{k}(s)[\mathrm T_1,\ldots,\mathrm T_m]$ engendré par les polynômes
$(P_i)_{(\boldsymbol{k}(s))}$ ($i\in I$). Alors, si
$V(\mathfrak a_\eta)=\varnothing$, il existe un voisinage $U$ de $\eta$ dans
$S$ tel que $V(\mathfrak a_s)=\varnothing$ pour tout $s\in U$.
\end{lemma}

En effet, soit $Y=\operatorname{Spec}(B)$, et soit $Z$ la partie fermée
$V(\mathfrak a)$ de $Y$~; comme $\mathfrak a$ est un idéal de type fini, $Z$
est constructible dans $Y$ $(\mathbf 0_{\mathrm{III}},9.1.5)$ et avec les
notations introduites en (9.4.1), on a $V(\mathfrak a_s)=Z_s$ pour tout
$s\in S$~; comme le morphisme structural $Y\to S$ est de présentation finie,
la conclusion du lemme résulte de (9.5.1) et (9.2.3).

